% Options for packages loaded elsewhere
\PassOptionsToPackage{unicode}{hyperref}
\PassOptionsToPackage{hyphens}{url}
\PassOptionsToPackage{dvipsnames,svgnames,x11names}{xcolor}
%
\documentclass[
  10pt,
]{article}
\usepackage{amsmath,amssymb}
\usepackage{iftex}
\ifPDFTeX
  \usepackage[T1]{fontenc}
  \usepackage[utf8]{inputenc}
  \usepackage{textcomp} % provide euro and other symbols
\else % if luatex or xetex
  \usepackage{unicode-math} % this also loads fontspec
  \defaultfontfeatures{Scale=MatchLowercase}
  \defaultfontfeatures[\rmfamily]{Ligatures=TeX,Scale=1}
\fi
\usepackage{lmodern}
\ifPDFTeX\else
  % xetex/luatex font selection
  \setmainfont[]{DejaVu Serif}
  \setsansfont[]{DejaVu Sans}
  \setmonofont[]{DejaVu Sans Mono}
  \setmathfont[]{Latin Modern Math}
\fi
% Use upquote if available, for straight quotes in verbatim environments
\IfFileExists{upquote.sty}{\usepackage{upquote}}{}
\IfFileExists{microtype.sty}{% use microtype if available
  \usepackage[]{microtype}
  \UseMicrotypeSet[protrusion]{basicmath} % disable protrusion for tt fonts
}{}
\makeatletter
\@ifundefined{KOMAClassName}{% if non-KOMA class
  \IfFileExists{parskip.sty}{%
    \usepackage{parskip}
  }{% else
    \setlength{\parindent}{0pt}
    \setlength{\parskip}{6pt plus 2pt minus 1pt}}
}{% if KOMA class
  \KOMAoptions{parskip=half}}
\makeatother
\usepackage{xcolor}
\usepackage[margin=2.0cm]{geometry}
\usepackage{longtable,booktabs,array}
\usepackage{calc} % for calculating minipage widths
% Correct order of tables after \paragraph or \subparagraph
\usepackage{etoolbox}
\makeatletter
\patchcmd\longtable{\par}{\if@noskipsec\mbox{}\fi\par}{}{}
\makeatother
% Allow footnotes in longtable head/foot
\IfFileExists{footnotehyper.sty}{\usepackage{footnotehyper}}{\usepackage{footnote}}
\makesavenoteenv{longtable}
\setlength{\emergencystretch}{3em} % prevent overfull lines
\providecommand{\tightlist}{%
  \setlength{\itemsep}{0pt}\setlength{\parskip}{0pt}}
\setcounter{secnumdepth}{-\maxdimen} % remove section numbering
% doc_header.tex -- LaTeX preamble additions used by build_docs.sh for every
% Markdown -> LaTeX -> PDF conversion (supplements, README, submission files).
\usepackage{etoolbox}
% longtable may not be loaded by pandoc when a document has no table
\usepackage{longtable}
% bold vectors: \bm -> \symbf under unicode-math (XeLaTeX), bm package otherwise
\makeatletter
\@ifpackageloaded{unicode-math}{\providecommand{\bm}[1]{\symbf{#1}}}{\usepackage{bm}}
\makeatother
% smaller type in long tables so that wide audit tables fit the text width;
% inside a table a line break is allowed after an underscore of plain text
% (pandoc writes it as \_), so that long identifiers wrap inside their column
% rows are set a little apart, so that symbols with indices do not touch
\AtBeginEnvironment{longtable}{\footnotesize\renewcommand{\arraystretch}{1.45}%
  \let\sptunderscore\_\def\_{\sptunderscore\allowbreak}}
\setlength{\LTcapwidth}{\textwidth}
% allow line breaks after slashes and hyphens inside table cells
\usepackage{xurl}
% code blocks: break over-long lines instead of running into the margin.
% Plain ``` blocks are typeset by pandoc as {verbatim}, highlighted ones
% (```bash) as {Highlighting}; both are re-declared as breakable fancyvrb
% environments.  The continuation line is marked with a small hooked arrow.
\usepackage{fvextra}
\RecustomVerbatimEnvironment{verbatim}{Verbatim}{breaklines,breakanywhere,fontsize=\small}
\makeatletter
\@ifundefined{Highlighting}{}{%
  \RecustomVerbatimEnvironment{Highlighting}{Verbatim}{breaklines,breakanywhere,fontsize=\small,commandchars=\\\{\}}}
\makeatother
% inline code (pandoc writes it as \texttt{...}): never hyphenate -- an inserted
% hyphen could be read as part of a file name -- but allow a line break after
% an underscore, slash, comma, colon or equals sign, so that long paths and
% column lists wrap inside the text width and inside table columns.
% Inside a formula (e.g. the name of a notebook symbol in an equation) the same
% command must first leave math mode, which \text of amsmath does.
\ExplSyntaxOn
\tl_new:N \l__spt_code_tl
\cs_new_protected:Npn \__spt_texttt:n #1
  {
    \tl_set:Nn \l__spt_code_tl {#1}
    \regex_replace_all:nnN { ([\/\,\:\=]) } { \1 \c{allowbreak} } \l__spt_code_tl
    \regex_replace_all:nnN { \c{_} } { \c{_} \c{allowbreak} } \l__spt_code_tl
    \group_begin: \ttfamily \hyphenchar \font = -1 \scan_stop: \l__spt_code_tl \group_end:
  }
\RenewDocumentCommand \texttt { m }
  {
    \mode_if_math:TF { \text { \__spt_texttt:n {#1} } } { \__spt_texttt:n {#1} }
  }
\ExplSyntaxOff
% A figure is written in the Markdown source as an image paragraph followed by
% a caption paragraph ("**Figure S1.** ..."), so that the caption is also
% visible on GitHub.  An image that starts a paragraph is centered.
\usepackage{graphicx}% pandoc loads it only for documents that contain an image
\NewCommandCopy\sptincludegraphics\includegraphics
\RenewDocumentCommand\includegraphics{O{} m}{%
  \ifvmode\begin{center}\sptincludegraphics[#1]{#2}\end{center}%
  \else\sptincludegraphics[#1]{#2}\fi}
% keep figures where they are referenced
\usepackage{float}
\floatplacement{figure}{H}
% running header with the document title
\usepackage{fancyhdr}
\pagestyle{fancy}
\fancyhf{}
\fancyhead[L]{\small PROJECT spacetimes --- documentation and provenance}
\fancyhead[R]{\small\thepage}
\renewcommand{\headrulewidth}{0.4pt}
% 16x16 matrices: tighter columns so that a full matrix fits the text width
\setlength{\arraycolsep}{3.2pt}
% The documents of this project contain long tagged symbols, file paths and
% identifiers that cannot be hyphenated.  A paragraph is therefore allowed to
% have looser lines before TeX gives up and lets a line stick out of the margin.
\AtBeginDocument{\tolerance=2000 \emergencystretch=5em \hbadness=2000}
\ifLuaTeX
  \usepackage{selnolig}  % disable illegal ligatures
\fi
\IfFileExists{bookmark.sty}{\usepackage{bookmark}}{\usepackage{hyperref}}
\IfFileExists{xurl.sty}{\usepackage{xurl}}{} % add URL line breaks if available
\urlstyle{same}
\hypersetup{
  pdftitle={P13 --- The gauge-covariant derivative for the spin groups of the 4+4 and 7+1 spacetimes, and Section 4 of Lord (1971)},
  pdfauthor={Prepared for Patrick Lee Nash, Ph.D.},
  colorlinks=true,
  linkcolor={blue},
  filecolor={Maroon},
  citecolor={blue},
  urlcolor={blue},
  pdfcreator={LaTeX via pandoc}}

\title{P13 --- The gauge-covariant derivative for the spin groups of the
4+4 and 7+1 spacetimes, and Section 4 of Lord (1971)}
\usepackage{etoolbox}
\makeatletter
\providecommand{\subtitle}[1]{% add subtitle to \maketitle
  \apptocmd{\@title}{\par {\large #1 \par}}{}{}
}
\makeatother
\subtitle{Provenance of PROJECT spacetimes: check and correction of the
author's summary of the Yang-Mills construction}
\author{Prepared for Patrick Lee Nash, Ph.D.}
\date{November 15, 2026}

\begin{document}
\maketitle

\hypertarget{what-this-document-is}{%
\section{1. What this document is}\label{what-this-document-is}}

On November 15, 2026, the author asked (i) to read and understand
Section 4, ``Change of representation as a gauge group'', of the paper
by E. A. Lord, \emph{General relativity from gauge invariance}, Proc.
Camb. Phil. Soc. \textbf{69}, 423 (1971), and (ii) to check, verify and
correct, in a provenance, his summary of the construction of Yang and
Mills. Sec. 2 states what Lord's Section 4 does. Sec. 3 quotes the
summary. Sec. 4 checks it step by step and lists the corrections. Sec. 5
gives the corrected construction in full for the spin groups of the
project. Every step is verified by the program (Sec. 6), whose complete
output is printed in Sec. 7. Every symbol is defined here; Sec. 8
collects them.

\hypertarget{section-4-of-lord-1971-in-brief}{%
\section{2. Section 4 of Lord (1971), in
brief}\label{section-4-of-lord-1971-in-brief}}

\emph{Lord's starting point} (his Sec. 1). On the four-dimensional
spacetime of signature (1,3), the free field equations for particles of
any spin, written with Dirac matrices in the Weyl form, are invariant
not only under Lorentz transformations but also under \emph{changes of
representation} of the Dirac matrices: \(\gamma\to S\gamma S^{-1}\)
together with a transformation of the spinor indices of the field by the
same matrix \(S\). Lord keeps the changes that preserve the Weyl form:
block-diagonal matrices whose upper block is a unimodular complex
\(2\times2\) matrix and whose lower block is the inverse of its
Hermitian conjugate. They change the Pauli matrices \(\sigma^\mu\), of
which the Dirac matrices are built, into \(S\sigma^\mu S^\dagger\) (his
Eqs. 1.21, 1.22).

\emph{Section 4.} The central idea of the paper is to let that matrix
depend on the point of spacetime, as in the theory of Yang and Mills.
Then:

\begin{enumerate}
\def\labelenumi{\arabic{enumi}.}
\tightlist
\item
  The Pauli matrices become four arbitrary, linearly independent
  Hermitian matrix fields \(\sigma^a(x)\); the Dirac matrices built from
  them, \(\gamma^a(x)\), transform as \(S(x)\gamma^a(x)S^{-1}(x)\) (his
  Eqs. 4.1 to 4.5).
\item
  The derivative \(\partial_\mu\) of the free equations is replaced by
  differential operators \(d_a\) that obey the rule of Leibniz but need
  not commute, \([d_a,d_b]=X_{ab}{}^c d_c\) (his Eqs. 4.6, 4.7).
\item
  Because the Lorentz transformation of a spinor is now built from the
  point-dependent \(\sigma^a(x)\), it depends on the point although its
  parameters are constants. The equations stay invariant only if \(d_a\)
  is replaced by a derivative \(D_a\) that subtracts auxiliary fields,
  one matrix for each index \(a\), acting on each spinor index of the
  field (his Eq. 4.13). The auxiliary fields have an inhomogeneous
  transformation law (his Eq. 4.14), and they are combinations of the
  generators \(\tfrac12(\gamma_a\gamma_b-\gamma_b\gamma_a)\) with
  coefficients \(a_{abc}\) and a universal coupling constant \(g\) (his
  Eqs. 4.15, 4.16).
\item
  With these auxiliary fields the generalized equations are invariant
  under the point-dependent changes of representation \(S(x)\) (his Eqs.
  4.30 to 4.32): the gauge group of the paper.
\item
  Reference vectors (a tetrad), the metric built from them, and the
  connection that makes their derivative vanish turn the construction
  into a theory with an affine connection; the commutator of two
  derivatives gives the curvature (his Eqs. 4.33 to 4.48). In later
  sections the auxiliary fields are identified with the potentials of a
  self-interacting spin-2 field, and the resulting theory with
  Einstein's.
\end{enumerate}

\emph{A remark on Lord's transformation law.} In Lord's notation the
inhomogeneous term of the transformation of the auxiliary fields is
written \(S\,d_aS^{-1}\) (his Eqs. 4.14, 4.31). With his derivative
\(D_a=d_a-(\text{auxiliary field})\) (his Eq. 4.13), the derivation of
Sec. 4 below shows that this term must equal \((d_aS)\,S^{-1}\). The
other possible reading, \(S\,(d_aS^{-1})=-(d_aS)\,S^{-1}\), has the
opposite sign and would not make \(D_a\) covariant; check \texttt{G3c}
of Sec. 7 verifies that the opposite sign fails. Lord's expression for
the curvature,
\(\partial_\nu\Gamma_\mu-\partial_\mu\Gamma_\nu+[\Gamma_\mu,\Gamma_\nu]\)
in his notation (his Eq. 4.48), is the sign that follows from
\(D_a=d_a-(\text{auxiliary field})\); check \texttt{G4} verifies it in
the notation of Sec. 5.

\emph{Pages read}: the pages 424 to 426 and 432 to 437 of the journal
(pages 2 to 4 and 10 to 15 of the file). The other pages were not read.

\hypertarget{the-summary-that-is-to-be-checked}{%
\section{3. The summary that is to be
checked}\label{the-summary-that-is-to-be-checked}}

The author wrote:

\begin{quote}
Recall Yang-Mills features: Let X = coordinates. D = partial derivative
with respect to X. Let T\_\{ij\} generate a (usually compact, but not
here) group S; (indices i and j). Consider field a Y that transforms
under S: Y -\textgreater{} Y' = S . Y ; Gauge this: Let Y'(X) =
exp(g*w(X)\^{}\{ij\} * T\_\{ij\}/2 ) . Y(X) (sum over repeated indices i
and j implied), \textbackslash def S(X) . Y(X), where S(X) =
exp(g*w(X)\^{}\{ij\} * T\_\{ij\}) , where a sum over repeated indices i
and j is implied. Then D Y'(X) = S(X) . D Y(X) + {[} D S(X) {]} . Y(X).
Define gauge field A: D S(X) /def A'(X) . S(X) - S(X) . A(X); Therefore
D Y'(X) = S(X) . D Y(X) + {[} A'(X) . S(X) - S(X) . A(X) {]} . Y(X);
thus `gauge covariant derivative' is D Y'(X) - A'(X) . Y'(X) = S(X) .
{[}D Y(X) - A(X) . Y(X) {]}
\end{quote}

\hypertarget{the-check-step-by-step}{%
\section{4. The check, step by step}\label{the-check-step-by-step}}

\textbf{The verdict.} The chain of equations is correct, and its
conclusion, the covariance of the derivative \(\partial-\mathcal A\), is
right. Five points need correction or completion.

\textbf{(C1) The factor 1/2.} The transformed field is defined with the
exponent \(g\,w^{ij}T_{ij}/2\), but \(S(X)\) is then written with the
exponent \(g\,w^{ij}T_{ij}\). The two differ by the factor 2, and the
derivation needs one and the same \(S\) in both places: with
\(Y'=\exp(\tfrac12gwT)Y\) and the gauge field transformed with
\(\exp(gwT)\), the final identity is false (check \texttt{G3d}: the two
sides differ by 71 and 106 per cent on the two spacetimes). The factor
\(\tfrac12\) belongs in both, because the parameters and the generators
are antisymmetric, \(w^{ij}=-w^{ji}\) and \(T_{ij}=-T_{ji}\), so that
the sum over all \(i,j\) counts every pair twice:
\(\tfrac12\sum_{i,j}w^{ij}T_{ij}=\sum_{i<j}w^{ij}T_{ij}\). For the spin
group, with \(T_{ij}\) replaced by
\(\Sigma^{AB}=\tfrac14[\gamma^A,\gamma^B]\), the element
\(S=\exp(\tfrac12\sum_{A,B}w_{AB}\Sigma^{AB})\) is exactly the spinor
image of the Lorentz transformation \(\exp(w\eta)\) of the vectors,
\(S\gamma^CS^{-1}=\gamma^A[\exp(w\eta)]_A{}^C\); without the
\(\tfrac12\) the angle is doubled (check \texttt{G2}).

\textbf{(C2) One symbol for two objects.} The letter \(S\) names the
group and also its element. Here the group is \(\mathrm{Spin}_0(4,4)\)
or \(\mathrm{Spin}_0(7,1)\), and \(S(X)\) its element.

\textbf{(C3) The derivative and the gauge field carry an index.} There
is one partial derivative for each of the eight coordinates,
\(\partial_C=\partial/\partial x^C\), and therefore one gauge field for
each, \(\mathcal A_C(X)\): a matrix field with values in the Lie algebra
of the group,
\(\mathcal A_C=\tfrac12g_{\rm YM}\sum_{A,B}\mathcal A_C{}^{AB}\Sigma_{AB}\).

\textbf{(C4) The gauge field is not defined by the equation; the
equation is its transformation law.} ``\(DS\equiv A'S-SA\)'' does not
define the gauge field: \(\mathcal A_C\) is a new field, with its own
values at every point. What the equation fixes is how it changes when
the field \(\psi\) is transformed with \(S(X)\):
\[\mathcal A'_C=S\,\mathcal A_C\,S^{-1}+(\partial_CS)\,S^{-1}.\] This
law is \emph{chosen} so that the derivative becomes covariant; it keeps
\(\mathcal A_C\) in the Lie algebra, because \(S\mathcal A_CS^{-1}\) and
\((\partial_CS)S^{-1}\) both lie in it (check \texttt{G3a}). For an
element close to the identity,
\(S\approx1+\tfrac12g_{\rm YM}w_{AB}\Sigma^{AB}\), it reads
\(\delta\mathcal A_C=\tfrac12g_{\rm YM}(\partial_Cw_{AB})\Sigma^{AB}+[\tfrac12g_{\rm YM}w_{AB}\Sigma^{AB},\mathcal A_C]\).
The sign of the inhomogeneous term matters: with
\(-(\partial_CS)S^{-1}\) the covariance fails (check \texttt{G3c}).

\textbf{(C5) The group is not compact, and that has two consequences
that the summary does not mention.} (a) \(S(X)\) is not unitary. A
generator \(\Sigma^{ab}=\tfrac12\gamma^a\gamma^b\) that mixes a
space-like with a time-like axis has the square
\((\Sigma^{ab})^2=+\tfrac14\), so that
\(\exp(\theta\,\Sigma^{ab})=\cosh\tfrac\theta2+2\Sigma^{ab}\sinh\tfrac\theta2\)
grows without bound with the real parameter \(\theta\); no unitary
matrix does. Hence \(\psi^\dagger\psi\) is not invariant, and invariant
expressions must be built with an invariant bilinear form, such as
\(\bar\psi\psi'\) with \(\bar\psi=\psi^\dagger\beta\) or
\(\psi^TC^-\psi'\). (b) The invariant form on the Lie algebra,
\(\mathrm{tr}(\Sigma^{AB}\Sigma^{CD})\), is indefinite: it is diagonal,
with \(\mathrm{tr}(\Sigma^{AB}\Sigma^{AB})=-4\eta^{AA}\eta^{BB}\),
positive for the 16 (on 4+4) or 7 (on 7+1) generators that mix the two
kinds of axes and negative for the 12 or 21 others (check \texttt{G1}).
A Lagrangian of the Yang-Mills form built with it, proportional to
\(\mathrm{tr}(\mathcal F_{CD}\mathcal F^{CD})\), gives kinetic energies
of both signs to the components of the gauge field. That is the known
obstacle to a Yang-Mills theory of a non-compact group, and a theory
that uses this construction must say how it avoids it. (Lord does not
give his auxiliary fields a Yang-Mills Lagrangian of their own.)

\textbf{What is right and stays.} The product rule
\(\partial_C(S\psi)=S\,\partial_C\psi+(\partial_CS)\,\psi\); the
substitution of the transformation law; and the conclusion
\(\partial_C\psi'-\mathcal A'_C\psi'=S\,(\partial_C\psi-\mathcal A_C\psi)\),
checked at random points with random fields to \(3\times10^{-13}\)
(check \texttt{G3b}).

\hypertarget{the-corrected-construction}{%
\section{5. The corrected
construction}\label{the-corrected-construction}}

Let \(\psi(X)\) be the sixteen-component field of the project,
\(X=\{x0,\dots,x7\}\) the coordinates,
\(\partial_C=\partial/\partial x^C\), and
\(\Sigma^{AB}=\tfrac14[\gamma^A,\gamma^B]\) the 28 generators of the
spin group in the spinor representation. Let \(w_{AB}(X)=-w_{BA}(X)\) be
28 real functions and \(g_{\rm YM}\) a constant.

\begin{enumerate}
\def\labelenumi{\arabic{enumi}.}
\tightlist
\item
  \emph{The transformation.}
  \[S(X)=\exp\Big(\tfrac12g_{\rm YM}\sum_{A,B}w_{AB}(X)\,\Sigma^{AB}\Big),\qquad\psi'(X)=S(X)\,\psi(X).\]
\item
  \emph{The product rule.}
  \(\partial_C\psi'=S\,\partial_C\psi+(\partial_CS)\,\psi\): the
  derivative of \(\psi\) does not transform like \(\psi\).
\item
  \emph{The gauge field.} Introduce eight matrix fields
  \(\mathcal A_C(X)=\tfrac12g_{\rm YM}\sum_{A,B}\mathcal A_C{}^{AB}(X)\,\Sigma_{AB}\)
  with values in the Lie algebra, and require that, when \(\psi\) is
  transformed with \(S\), they change into
  \[\mathcal A'_C=S\,\mathcal A_C\,S^{-1}+(\partial_CS)\,S^{-1},\qquad\text{equivalently}\qquad\partial_CS=\mathcal A'_CS-S\,\mathcal A_C .\]
\item
  \emph{The covariant derivative}
  \(\mathcal D_C=\partial_C-\mathcal A_C\) then satisfies
  \[\partial_C\psi'-\mathcal A'_C\psi'=S\,\partial_C\psi+(\mathcal A'_CS-S\mathcal A_C)\psi-\mathcal A'_CS\psi=S\,(\partial_C\psi-\mathcal A_C\psi),\]
  that is, \(\mathcal D'_C\psi'=S\,\mathcal D_C\psi\).
\item
  \emph{The field strength.}
  \(\mathcal F_{CD}=\partial_C\mathcal A_D-\partial_D\mathcal A_C-[\mathcal A_C,\mathcal A_D]\)
  transforms as \(\mathcal F'_{CD}=S\,\mathcal F_{CD}\,S^{-1}\), and
  \([\mathcal D_C,\mathcal D_D]\psi=-\mathcal F_{CD}\,\psi\) (check
  \texttt{G4}).
\item
  \emph{The chiral parts.} The generators commute with the chirality
  matrix \(\chi\) (equal to \(\gamma^{0}\cdots\gamma^{7}\) on 4+4 and
  \(i\gamma^{0}\cdots\gamma^{7}\) on 7+1), and so do \(S\) and
  \(\mathcal A_C\). The construction therefore acts within the
  left-handed and within the right-handed spinors, \(P^{L,R}\psi\) with
  \(P^{L,R}=\tfrac12(1\mp\chi)\), with the \emph{same} parameters
  \(w_{AB}\) and the same gauge field. Independent transformations of
  the two parts would need two independent sets of parameters and two
  gauge fields: a larger group, which is a choice of the model and not
  part of this construction.
\end{enumerate}

\hypertarget{what-was-computed-how-to-reproduce-it-and-what-was-not-computed}{%
\section{6. What was computed, how to reproduce it, and what was not
computed}\label{what-was-computed-how-to-reproduce-it-and-what-was-not-computed}}

The command \texttt{gaugecov} of the program of the main project (file
\texttt{code/spacetimes-rs/src/gaugecov\_cmd.rs}) makes the checks
printed in Sec. 7, for the 4+4 and for the 7+1 spacetime: the trace form
in exact integers (\texttt{G1}); the relation of \(S\) to the Lorentz
transformation for random parameters (\texttt{G2}); and, at four random
points of the plane of the coordinates \(x1\) and \(x4\), with smooth
random functions \(w_{AB}\), \(\mathcal A_1\), \(\mathcal A_4\) and
\(\psi\), the transformation of the gauge field, the covariance, the two
failures, the field strength and the commutator (\texttt{G3a} to
\texttt{G4}). Derivatives are central differences of fourth order with
the step \(10^{-3}\), and the exponential of a matrix is computed by
scaling and squaring with a Taylor polynomial of degree 18. From the
root of the repository:

\begin{verbatim}
cargo run --release --manifest-path code/spacetimes-rs/Cargo.toml -- gaugecov
\end{verbatim}

It writes \texttt{data/gaugecov\_checks.txt} and ends with ``gaugecov: N
of N checks passed''. This document is built with

\begin{verbatim}
python3 code/scripts/assemble_docs.py docs_src/P13_gauge_covariant_derivative.src.md
bash code/scripts/build_docs.sh provenance/P13_gauge_covariant_derivative.md
\end{verbatim}

\emph{Not computed:} a Lagrangian for the gauge field; Lord's
identification of the auxiliary fields with a spin-2 field and with
Einstein's theory (his later sections); the pages of the paper listed as
not read in Sec. 2. The documentation of the program
(\texttt{provenance/P8\_code}) does not yet list the command
\texttt{gaugecov}.

\emph{Input:} the file
\texttt{General\ relativity\ from\ gauge\ invariance.pdf} (Lord 1971, 20
pages, a scan with recognized text); the file is not in the repository.
Its SHA-256 checksum:

\begin{verbatim}
c928b11e5721bb37093f9b19b19e8e9d4a743fcd3f8f4779e804845fcaff2275
\end{verbatim}

\hypertarget{the-complete-output-of-the-checks}{%
\section{7. The complete output of the
checks}\label{the-complete-output-of-the-checks}}

\begin{verbatim}
# PROJECT spacetimes: the gauge-covariant derivative for the spin groups of the 4+4 and 7+1 spacetimes, checked numerically
# generated by: spacetimes gaugecov (code/spacetimes-rs/src/gaugecov_cmd.rs); trace form in exact integers, the rest in double precision
# Sigma^{AB} = (1/4)[gamma^A, gamma^B]; S = exp((1/2) g_YM w_{AB} Sigma^{AB}); cal-D_C = d_C - cal-A_C; cal-A'_C = S cal-A_C S^-1 + (d_C S) S^-1
[PASS] 44-G1  4+4: the trace form tr(Sigma^{AB} Sigma^{CD}) of the 28 generators is diagonal, with tr(Sigma^{AB} Sigma^{AB}) = -4 eta^{AA} eta^{BB}: 16 positive (the generators that mix a space-like with a time-like axis, boosts) and 12 negative (rotations); the invariant form is indefinite because the group is not compact
[PASS] 44-G2  4+4: for random parameters w_{AB} = -w_{BA}, S = exp((1/2) w_{AB} Sigma^{AB}) (sum over all A and B) satisfies S gamma^C S^-1 = gamma^A [exp(w eta)]_A^C to 2.4e-15 (relative): S is the spinor image of the Lorentz transformation exp(w eta). With exp(w_{AB} Sigma^{AB}), the exponent without the factor 1/2, one obtains exp(2 w eta) instead (to 2.3e-14); it differs from exp(w eta) by 2.20 (relative)
[PASS] 44-G3a  4+4: at four random points and for C = 1, 4, the transformed gauge field cal-A'_C = S cal-A_C S^-1 + (d_C S) S^-1 lies in the Lie algebra spanned by the 28 generators: the part outside it is at most 1.1e-12 of its size
[PASS] 44-G3b  4+4: with psi' = S psi, S = exp((1/2) g_YM w_{AB}(x) Sigma^{AB}), the identity d_C psi' - cal-A'_C psi' = S (d_C psi - cal-A_C psi) holds to 2.8e-13 (relative; the derivatives by central differences of fourth order with the step 1e-3)
[PASS] 44-G3c  4+4: with the opposite sign of the inhomogeneous term, cal-A'_C = S cal-A_C S^-1 - (d_C S) S^-1 = S cal-A_C S^-1 + S d_C(S^-1), the identity fails: the two sides differ by at least 0.40 (relative)
[PASS] 44-G3d  4+4: with psi' = exp((1/2) g_YM w Sigma) psi but the gauge field transformed with exp(g_YM w Sigma), as the two formulas of the recalled text would have it, the identity fails: the two sides differ by at least 0.71 (relative)
[PASS] 44-G4  4+4: the field strength cal-F_14 = d_1 cal-A_4 - d_4 cal-A_1 - [cal-A_1, cal-A_4] transforms as S cal-F_14 S^-1 (to 3.3e-11), and [cal-D_1, cal-D_4] psi = -cal-F_14 psi (to 3.9e-11), with cal-D_C = d_C - cal-A_C
[PASS] 71-G1  7+1: the trace form tr(Sigma^{AB} Sigma^{CD}) of the 28 generators is diagonal, with tr(Sigma^{AB} Sigma^{AB}) = -4 eta^{AA} eta^{BB}: 7 positive (the generators that mix a space-like with a time-like axis, boosts) and 21 negative (rotations); the invariant form is indefinite because the group is not compact
[PASS] 71-G2  7+1: for random parameters w_{AB} = -w_{BA}, S = exp((1/2) w_{AB} Sigma^{AB}) (sum over all A and B) satisfies S gamma^C S^-1 = gamma^A [exp(w eta)]_A^C to 2.5e-15 (relative): S is the spinor image of the Lorentz transformation exp(w eta). With exp(w_{AB} Sigma^{AB}), the exponent without the factor 1/2, one obtains exp(2 w eta) instead (to 1.3e-14); it differs from exp(w eta) by 2.06 (relative)
[PASS] 71-G3a  7+1: at four random points and for C = 1, 4, the transformed gauge field cal-A'_C = S cal-A_C S^-1 + (d_C S) S^-1 lies in the Lie algebra spanned by the 28 generators: the part outside it is at most 9.9e-13 of its size
[PASS] 71-G3b  7+1: with psi' = S psi, S = exp((1/2) g_YM w_{AB}(x) Sigma^{AB}), the identity d_C psi' - cal-A'_C psi' = S (d_C psi - cal-A_C psi) holds to 3.0e-13 (relative; the derivatives by central differences of fourth order with the step 1e-3)
[PASS] 71-G3c  7+1: with the opposite sign of the inhomogeneous term, cal-A'_C = S cal-A_C S^-1 - (d_C S) S^-1 = S cal-A_C S^-1 + S d_C(S^-1), the identity fails: the two sides differ by at least 0.28 (relative)
[PASS] 71-G3d  7+1: with psi' = exp((1/2) g_YM w Sigma) psi but the gauge field transformed with exp(g_YM w Sigma), as the two formulas of the recalled text would have it, the identity fails: the two sides differ by at least 1.06 (relative)
[PASS] 71-G4  7+1: the field strength cal-F_14 = d_1 cal-A_4 - d_4 cal-A_1 - [cal-A_1, cal-A_4] transforms as S cal-F_14 S^-1 (to 3.0e-11), and [cal-D_1, cal-D_4] psi = -cal-F_14 psi (to 2.8e-11), with cal-D_C = d_C - cal-A_C
14 of 14 checks passed
\end{verbatim}

\hypertarget{notation}{%
\section{8. Notation}\label{notation}}

\begin{longtable}[]{@{}
  >{\raggedright\arraybackslash}p{(\columnwidth - 2\tabcolsep) * \real{0.3649}}
  >{\raggedright\arraybackslash}p{(\columnwidth - 2\tabcolsep) * \real{0.6351}}@{}}
\toprule\noalign{}
\begin{minipage}[b]{\linewidth}\raggedright
symbol
\end{minipage} & \begin{minipage}[b]{\linewidth}\raggedright
meaning
\end{minipage} \\
\midrule\noalign{}
\endhead
\bottomrule\noalign{}
\endlastfoot
\(X=\{x0,\dots,x7\}\), \(x^C\) & the coordinates; the coordinate with
index \(C\) \\
\(\partial_C\) & \(\partial/\partial x^C\) \\
\(\gamma^A\), \(\eta^{AB}\), \(A,B,C,D\) & the real \(16\times16\) Dirac
matrices with \(\gamma^A\gamma^B+\gamma^B\gamma^A=2\eta^{AB}\); the flat
metric, \(\mathrm{diag}(+1,+1,+1,+1,-1,-1,-1,-1)\) (4+4) or
\(\mathrm{diag}(-1,-1,-1,-1,+1,-1,-1,-1)\) (7+1); indices
\(0,\dots,7\) \\
\(\Sigma^{AB}\), \(\Sigma_{AB}\) & \(\tfrac14[\gamma^A,\gamma^B]\), the
generators of the spin group; the same with lowered indices \\
\(\mathrm{Spin}_0(4,4)\), \(\mathrm{Spin}_0(7,1)\) & the identity
components of the spin groups \\
\(w_{AB}(X)\), \(\theta\) & the 28 real parameters of a point-dependent
element of the spin group; one real parameter (in C5 only) \\
\(g_{\rm YM}\) & the coupling constant of the gauge field \\
\(S(X)\), \(S^{-1}\) & the element
\(\exp(\tfrac12g_{\rm YM}w_{AB}\Sigma^{AB})\) of the spin group in the
spinor representation; its inverse \\
\(\exp(w\eta)\) & the Lorentz transformation of the vectors, the
exponential of the \(8\times8\) matrix with the entries
\(w_{AC}\eta^{CC}\) \\
\(\psi\), \(\psi'\) & the sixteen-component field; the transformed
field \\
\(\mathcal A_C\), \(\mathcal A_C{}^{AB}\), \(\mathcal A'_C\) & the gauge
field; its components; the transformed gauge field \\
\(\mathcal D_C\) & \(\partial_C-\mathcal A_C\), the covariant
derivative \\
\(\mathcal F_{CD}\) & the field strength \\
\(\chi\), \(P^L\), \(P^R\) & the chirality matrix; the chiral projectors
\(\tfrac12(1\mp\chi)\) \\
\(\bar\psi\), \(\beta\), \(C^-\) & \(\psi^\dagger\beta\) with
\(\beta=-i\gamma^{4}\) (4+4) or \(\gamma^{4}\) (7+1); the
charge-conjugation matrix, \(\gamma^{0}\gamma^{1}\gamma^{2}\gamma^{3}\)
(4+4) or \(\gamma^{4}\gamma^{0}\cdots\gamma^{7}\) (7+1) \\
\(Y\), \(T_{ij}\), \(w^{ij}\), \(g\), \(A\), \(D\) & the author's
symbols in the quotation of Sec. 3 only: the field, the generators, the
parameters, the coupling, the gauge field, the partial derivative;
elsewhere this document uses \(\psi\), \(\Sigma^{AB}\), \(w_{AB}\),
\(g_{\rm YM}\), \(\mathcal A_C\), \(\partial_C\) \\
\(\sigma^\mu\), \(\sigma^a(x)\), \(\gamma^a(x)\), \(d_a\),
\(X_{ab}{}^c\), \(a_{abc}\), \(g\), \(\Gamma_\mu\) & Lord's symbols in
Sec. 2 only: the Pauli matrices, their point-dependent version, the
Dirac matrices built from them, his derivatives, their commutator
coefficients, the coefficients and the coupling of his auxiliary fields,
the auxiliary fields \\
\end{longtable}

\end{document}
